\chapter{Tier 1: deterministic investment appraisal}\label{chap:tier1}
\section{Problem analysis}\label{sec:problem1}
We select battery and electrolyser capacities and their daily use over
2020--25 (2,192 consecutive days). The investment is made once and applies
to all six years. Direct electricity sales provide the no-investment
benchmark. The investment case depends on three questions: whether
electrolysis can use electricity that could otherwise be sold, how much
hydrogen a buyer can take, and whether storage can shift enough electricity
to recover its costs.

\section{Assumptions and research}\label{sec:research1}
\subsection{Tier-specific assumptions}
\emph{Policy A} allocates electricity by economic value.
\emph{Policy B} serves current electricity demand first with generation
and then with stored electricity. B honours this priority even at negative
prices. Hydrogen is sold as produced subject to a common daily offtake limit.
Battery electricity can feed electrolysis after any applicable sales
priority. The daily model carries battery inventory forward, with at most
one net charge or discharge per model per day and zero opening and closing usable
inventory over the study period. Daily operation uses the historical paths.

We distinguish a broad market case and two smaller buyer cases.
The broad cap of 1,559 t/day is the highest UK annual hydrogen consumption
estimate in the 2022--24 source divided by the number of days in that year.
The 1.56 and 10 t/day cases illustrate smaller outlets.
Netback is the revenue per kg after non-electricity variable costs.
We take \gbp 7.80/kg as a scenario, rebasing the UK policy benchmark of
\gbp 175/MWh in 2012 pounds using 39.4 kWh/kg and the consumer price index.
The national total includes existing uses and provides a scale comparison
for the constructed system \cite{har,eho,cpi,basis}.

\subsection{Research objectives and parameter selection}
\begin{table}[H]
\centering\small
\caption{Physical and financial inputs. Monetary values are in 2020 pounds.}
\label{tab:research}
\begin{tabularx}{\linewidth}{@{}P{.19\linewidth}YY@{}}
\toprule
Input & Model specification & Basis and use \\
\midrule
Battery options &
Four Sungrow, Huawei and Tesla models, 0.215--3.854 MWh and
0.108--1.927 MW per unit. Round-trip efficiency is 90--92\%,
state of charge is limited to 20--80\%, and life is 15 years. &
Supplier data for performance. Course and Danish Energy Agency costs for
the smaller units and a public Tesla installed quote for the large unit.
Capital is \gbp 55,129--1,239,440 per unit and annual fixed operation is
\gbp 1,599--30,986 \cite{brief,dea,v6}. \\
Electrolyser options &
Eight ITM and Nel models, 1.31--50.31 MW. Specific consumption is
55.9--59.0 kWh/kg and the minimum online load is 10--25\%, by model. &
Manufacturer operating data. Common installed cost of \gbp 1.46m/MW
lets us compare size and efficiency on the same financial basis \cite{iea,v6}. \\
Electrolyser upkeep &
\gbp 25,862/MW-year fixed operation, \gbp 489,609/MW per stack renewal in
years 11 and 22, 30-year life. &
UK cost study. A second case uses its complete cost basis with a reserve
per kg of output in place of scheduled stack renewal \cite{beis,v6}. \\
Financial treatment &
5\% real discount rate, zero residual value and battery commissioning charge. &
Costs enter as annual equivalents. The 20\% starting state of charge is
accounted for through commissioning \cite{basis,v6}. \\
\bottomrule
\end{tabularx}
\end{table}

\newpage
\section{Mathematical model}\label{sec:model1}
Indices \(i\) and \(j\) denote battery and electrolyser models.
All daily energy flows and inventories are in MWh. Installed counts
\(b_i,n_j\) and online electrolyser counts \(z_{jt}\) are nonnegative
integers. The daily variables are direct sales \(q_t\), generation sent
directly to electrolysis \(g_t^H\), battery charge \(c_{it}\),
discharge for sales \(r_{it}\) and for electrolysis \(u_{it}\), usable stock
\(s_{it}\), and electrolyser input \(h_{jt}\). The conversion factor
\(\kappa_j\) is MWh/kg and \(L\) is the daily offtake limit in kg.
With netback \(v\) in \gbp/kg and annual equipment costs \(a_i^B,a_j^H\),
we maximise annual equivalent profit
\begin{equation}
\max\quad \Pi =
\frac{1}{Y}\sum_{t\in\mathcal T}
\left[p_t\left(q_t+\sum_i r_{it}\right)+
\sum_j (v-w_j)\frac{h_{jt}}{\kappa_j}\right]
-\sum_i a_i^B b_i-\sum_j a_j^H n_j .
\label{eq:profit}
\end{equation}
Here \(w_j\) is the variable stack reserve in \gbp/kg, zero when
stack renewal is charged through the annual equipment cost.

The daily flow equations allocate generation, cap total sales at market
demand and balance both sources of electricity into electrolysis:
\begin{align}
 q_t+g_t^H+\sum_i c_{it}&\le G_t, &
 q_t+\sum_i r_{it}&\le D_t, &
 \sum_j h_{jt}&=g_t^H+\sum_i u_{it}. \label{eq:flow}
\end{align}
The unallocated part of generation is unused. For battery \(i\), \(W_i\)
is the usable energy per unit, \(P_i\) the power in MW, and \(\eta_i\)
the square root of round-trip efficiency. The inventory and daily throughput
constraints are
\begin{align}
s_{it}&=(1-\delta_i)s_{i,t-1}+\eta_i c_{it}-(r_{it}+u_{it})/\eta_i,
 & 0\le s_{it}&\le W_i b_i, \label{eq:stock}\\
c_{it}&\le b_i\min(24P_i,W_i/\eta_i),
 & r_{it}+u_{it}&\le b_i\min(24P_i,\eta_i W_i). \label{eq:window}
\end{align}
Daily self-discharge \(\delta_i\) is zero in the main case. We set
\(s_{i0}=s_{iT}=0\) for \(T=|\mathcal T|\), and at least one of \(c_{it}\)
and \(r_{it}+u_{it}\) is zero. A binary charge-mode indicator enforces
this choice when needed. For electrolysis, \(H_j\) is the full-load daily
input per unit and \(\alpha_j\) its minimum operating fraction:
\begin{align}
\alpha_j H_j z_{jt}\le h_{jt}&\le H_j z_{jt}, &
z_{jt}&\le n_j, &
\sum_j h_{jt}/\kappa_j&\le L. \label{eq:hydrogen}
\end{align}
All flows are nonnegative. With \((x)_+=\max(0,x)\), B also imposes
\begin{equation}
q_t=\min(G_t,D_t), \qquad
\sum_i r_{it}=\min\!\left\{(D_t-G_t)_+,\,
            \sum_i\eta_i(1-\delta_i)s_{i,t-1}\right\}.
\label{eq:priority}
\end{equation}
The short-duration battery models can discharge all usable stock within
a day. A binary indicator selects the binding branch of the minimum.

For a unit with capital \(K\), fixed annual operating cost \(F\), lifetime
\(N\), real discount rate \(d\) and renewals \(C_k\) in years \(k\),
annual cost is
\begin{equation}
a=F+\frac{d}{1-(1+d)^{-N}}
\left(K+\sum_k\frac{C_k}{(1+d)^k}\right).
\label{eq:annualcost}
\end{equation}
Reported gain subtracts each policy's own optimum with all \(b_i,n_j=0\).

\textbf{Numerical checks.}
HiGHS certified the 14 portfolio cases with a maximum reported relative
gap below \(10^{-9}\). Independent checks reproduced balances, cash
flows and zero-battery dispatch. Six forced-storage cases tested stock
carry-over, discharge to electrolysis and B's priority rule \cite{v6}.

\newpage
\section{Results and analysis}\label{sec:results1}
\subsection{Capacity and historical performance}
In the broad market case, electrolysis-only and joint investment return
the same portfolios. A selects 3,631 MW (721 Neptune V and two Neptune II
units) and B selects 740 MW (147 Neptune V units). At 10 t/day, both choose
23.4 MW across six smaller modules. Initial capital in the broad case is
\gbp 5.30bn (A) and \gbp 1.08bn (B), producing about 488,000 and
19,100 t/year of hydrogen respectively.
\begin{table}[H]
\centering\small
\caption{Six-year portfolio choices and annual profit gains by policy.
A dash marks an excluded investment. Gains are in \gbp m/year.}
\label{tab:annual}
\textit{(a) Allowed assets at the broad market cap}\par\vspace{2pt}
\begin{tabular}{@{}l r r r r r r@{}}
\toprule
 & \multicolumn{3}{c}{Policy A} & \multicolumn{3}{c}{Policy B} \\
\cmidrule(lr){2-4}\cmidrule(l){5-7}
Allowed investment & Battery & Electrolysis & Gain & Battery & Electrolysis & Gain \\
 & (MWh) & (MW) & & (MWh) & (MW) & \\
\midrule
Battery only & 0 & -- & 0 & 0 & -- & 0 \\
Electrolysis only & -- & 3,631 & 1,639 & -- & 740 & 37.71 \\
Joint choice & 0 & 3,631 & 1,639 & 0 & 740 & 37.71 \\
\bottomrule
\end{tabular}
\par\vspace{5pt}
\textit{(b) Smaller outlets with both assets allowed}\par\vspace{2pt}
\begin{tabular}{@{}l r r r r r@{}}
\toprule
Daily offtake & Battery & \multicolumn{2}{c}{Electrolysis (MW)} &
 \multicolumn{2}{c}{Gain} \\
\cmidrule(lr){3-4}\cmidrule(l){5-6}
(t/day) & (MWh), A/B & A & B & A & B \\
\midrule
1.56 & 0 & 3.84 & 3.84 & 1.91 & 0.35 \\
10 & 0 & 23.42 & 23.42 & 12.68 & 2.44 \\
\bottomrule
\end{tabular}
\end{table}

At 10 t/day, B loses money in 2021, 2023 and 2024, with a worst
year of \gbp 1.26m below no investment. These years have 43, 29 and
30 surplus days, respectively, against 154 in 2022. A can produce on
additional days when the hydrogen value exceeds
the price of forgone electricity sales. At the broad cap, A meets 82.2\%
of electricity demand and B meets 92.2\%. Using the alternative UK cost
basis at that cap leaves A's capacity unchanged and raises B's from
740 to 1,465 MW, with zero batteries in both cases \cite{beis,v6}.

\subsection{Energy attribution and contract sensitivity}\label{sec:contract1}
The historical mean physical surplus is 2,276 GWh/year.
For the zero-battery solutions we attribute \(\min(g_t^H,S_t)\) to physical
surplus first. The remaining electrolysis input uses energy within current
electricity demand. We calculate the monetary loss separately against
same-policy electricity sales without investment, accounting for negative
prices under A. The aggregate records identify the available fuels by day
and cannot trace the generating unit that supplied electrolysis.
\begin{table}[H]
\centering\small
\caption{Electrolysis sourcing and forgone electricity receipts.
Energy is in GWh/year. The revenue reduction is in \gbp m/year.}
\label{tab:attribution}
\begin{tabular}{@{}l c r r r r@{}}
\toprule
Offtake & Policy & \shortstack{Total input} & \shortstack{Physical\\surplus} &
\shortstack{Within current\\demand} & \shortstack{Electricity\\revenue reduction} \\
\midrule
Broad & A & 27,255 & 2,184 & 25,070 & 1,618 \\
Broad & B & 1,067 & 1,067 & 0 & 0 \\
10 t/day & A & 183 & 43 & 140 & 9.18 \\
10 t/day & B & 43 & 43 & 0 & 0 \\
\bottomrule
\end{tabular}
\end{table}

In the broad market case, A uses about 2,184 GWh/year of physical surplus,
only 8\% of its total electrolysis input. Its electricity receipts fall by
\gbp 1,618m/year. B uses 1,067 GWh/year of surplus and preserves direct
sales, limiting its hydrogen throughput. At 10 t/day, each captures
about 43 GWh/year, just under 2\% of the physical surplus.

\newpage
\noindent
After reoptimising electrolysis and dispatch, the lowest netbacks that
support any investment are about \gbp 3.32/kg in A and \gbp 4.57/kg in B
at 10 t/day. At 1.56 t/day they rise to \gbp 3.38/kg and \gbp 4.66/kg.
The 23.4 MW portfolio in Table~\ref{tab:annual} is sized at \gbp 7.80/kg.
Holding its operation fixed, it breaks even at \gbp 3.90/kg in A or
\gbp 4.61/kg in B. At other prices, reoptimise operation and capacity.

\subsection{Storage and investment sensitivity}\label{sec:storage1}
\begin{table}[H]
\centering\small
\caption{Storage entry under historical and stressed conditions.
Price multipliers scale the standard deviation of daily prices.}
\label{tab:storage}
\begin{tabularx}{\linewidth}{@{}P{.35\linewidth}Y@{}}
\toprule
Test & Implication for storage investment \\
\midrule
Historical high-volatility windows &
All four battery options remain unselected in the highest-variance
60-, 90-, 180- and 365-day windows under both policies. The 90-day price
standard deviation is \gbp 114/MWh against \gbp 65/MWh over six years. \\
Battery cost and price stress &
Halving purchase and maintenance costs still gives zero batteries.
With historical supply and demand, batteries first enter at about
\(40\times\) (A) and \(127\times\) (B) the historical price standard
deviation at the modelled costs. \\
Joint supply and price stress &
The joint grid selects zero batteries at the historical price dispersion
across nine supply--demand gap dispersions. At \(120\times\) price
dispersion in B, raising gap dispersion from \(1\times\) to \(1.1\times\)
changes the optimum from zero to 2.60 GWh. \\
Cavern compressed air &
A separate 50 MW / 1.5 GWh test, at one-third of the reference battery
installation cost per MWh and 62\% round-trip efficiency, needs
\(5.38\)--\(5.50\times\) (A) or \(15.00\)--\(15.13\times\) (B) price
dispersion at the original supply--demand gap. \\
\bottomrule
\end{tabularx}
\end{table}
The volatility tests centre each series before scaling, preserving
means and date order. The battery and cavern tests allow the electrolyser
choice to adjust. For battery-only investment, even free charging and
the full usable discharge on every positive-price day yield at most
57\% of annual ownership cost across the four models. At historical
prices the cavern's annual operating
gain is \gbp 1.38m in A and \gbp 2.35m in B, against an annual cost of
\gbp 8.52m \cite{stress}.

\subsection{Investment recommendations}\label{sec:recommendations1}
\textbf{Size electrolysis to secured demand.}
For 10 t/day, appraise the 23.4 MW portfolio against the \gbp 34.2m
initial cost. A and B yield \gbp 12.68m and \gbp 2.44m/year respectively
(Table~\ref{tab:annual}). Confirm the ability to displace electricity
sales before relying on A. Under B, use the observed loss years to test
cash reserves and the case for staged procurement. Reduce the appraisal to roughly
3.84 MW if the outlet is near 1.56 t/day. At a netback below the
market entry levels above, defer electrolysis at the tested costs.
The broad cap describes market scale. Capacity at that scale warrants
appraisal only with much larger contracted offtake.

\textbf{Defer daily storage.}
Both cost and historical stress tests support zero batteries
(Table~\ref{tab:storage}). Reopen procurement if a lower installed cost
or additional revenue brings the expected operating gain above annual
ownership costs. For the cavern, the historical gain implies an annual
cost ceiling of \gbp 1.38m under A or \gbp 2.35m under B on these inputs.
